% ===================================================================
%  नकसा नं. ६ — घर के भीतर, सामान, खाना, अँजोर।
%
%  Traduction, faite en 2026, en bhojpouri d'aujourd'hui, sur l'ido
%  de texto/io. Regles de la colonne : voir 10-nakasa-01.tex. CINQ
%  pieces, cinq numerotations qui repartent de 1 (t06-c1-* a
%  t06-c5-*). 62 blocs, 211 renvois — le plus charge du livret —,
%  quarante-huit paires a retourner. DEUX notes d'auteur.
%
%  ECART DECLARE : t06-c1-06-1. L'ido y ouvre sur « (6) », qui est le
%  savon, la ou le francais porte « (16) », la femme de chambre.
%  renvoji.py laisse donc ce bloc passer SANS RIEN VERIFIER. La
%  colonne ecrit « (6) » comme les quarante-six autres — la source ne
%  bouge pas — et l'ordre a ete compte a la main : 6, 17, 20, 21, 22,
%  23, 24, 25, 26, 27, 32, 19, 18, treize appels.
%
%  LE PAIN (12) EST « डबलरोटी », ET C'EST UN REFUS A L'ENVERS. Les
%  quatre refus precedents de la colonne — चइत pour les mois, राम-राम
%  pour le salut, पतरा pour le calendrier, कबड्डी pour le jeu de
%  barres — ecartaient tous le mot le plus juste de la langue au
%  profit d'un mot neutre. Celui-ci fait le contraire : रोटी est le
%  mot ordinaire du pain, et il ne convient pas, parce que la रोटी de
%  la plaine est plate et cuite sur une plaque, tandis que le pain de
%  cette table a leve. La langue a deja fabrique le mot qu'il faut —
%  डबलरोटी, « le pain double » —, et il dit exactement ce qu'il faut
%  dire : c'est du pain, et ce n'est pas le notre. LA MEME LOI PRISE
%  PAR L'AUTRE BOUT : on ne transpose pas la chose, donc on prend le
%  mot marque et non le mot courant.
%
%  ET LE LIT (17) N'EST PAS UNE « खटिया ». Le lit de la plaine est un
%  cadre de bois sangle de corde, et il n'a ni sommier, ni matelas de
%  laine, ni traversin, ni edredon, ni ciel de lit — or l'alinea les
%  nomme tous les cinq. On ecrit « पलंग », qui est le lit a chassis.
%  Deux objets de mobilier dans le meme tableau, deux fois la meme
%  decision.
%
%  QUATRE AGES DE LA LUMIERE DANS UNE SEULE MAISON, ET LE VOCABULAIRE
%  SE DATE TOUT SEUL. Le titre du tableau annonce « la Lumizado », et
%  la maison en montre quatre etats : la bougie (31, 16) मोमबत्ती, la
%  lampe a petrole (12) मिट्टी तेल के लालटेन, la lampe electrique
%  (39) बिजली के बत्ती, le bec de gaz (62) गैस के मुँह. Le bhojpouri
%  n'a de mot a lui que pour le plus ancien, दिया, la lampe a huile —
%  et celui-la n'est pas dans le tableau. Tout le reste est
%  emprunte : मोम est persan, लालटेन est l'anglais lantern, गैस et
%  बत्ती electrique parlent d'eux-memes. UNE LANGUE NOMME PAR SON
%  FONDS CE QU'ELLE A TOUJOURS EU, ET PAR EMPRUNT CE QUI EST ARRIVE.
%
%  QUARANTE-HUIT PAIRES, QUATRE MANQUEES, ET LA COURBE SE REDRESSE.
%  Le tableau 3 en avait manque six sur vingt-deux, le tableau 5 cinq
%  sur trente-six, celui-ci quatre sur quarante-huit — soit un taux
%  qui tombe de vingt-sept a huit pour cent pendant que le nombre de
%  paires double. Les cinq chambres ont ete ecrites et relues
%  separement, comme cinq tableaux courts, et c'est la seule chose
%  qui ait change. LE REMEDE N'EST PAS DE MIEUX LIRE L'OUTIL, C'EST
%  DE COUPER LE TABLEAU EN MORCEAUX QUE LA VIGILANCE PEUT TENIR.
%  Les quatre restantes sont le parapluie (11) derriere son
%  porte-parapluies, le livre (19) derriere sa bibliotheque, le
%  tableau (37) derriere sa cloison et la coupe (10) derriere ses
%  fruits : toutes les quatre sont le meme locatif que le bhojpouri
%  pose devant, et aucune n'est ailleurs qu'en tete de piece.
% ===================================================================

%%K t06-apar-1 apar
\VUsaut{6.0mm}
\VUtitre{15.0pt}{60}{नकसा नं. ६}
\VUsaut{1.4mm}
\VUtitre{11.4pt}{0}{घर के भीतर, सामान, खाना,}
\VUsaut{0.4mm}
\VUtitre{11.4pt}{0}{अँजोर।}
\VUsaut{2.2mm}

%%K t06-apar-2 apar
घर बन गइल। योहानेस के चाचा आपन नया घर में रहे लागल बाड़ें, जेकर बनावट
हमनी के जानत बानी जा। अबहीं देखल जाव कि ऊ आपन घर के कइसे सजवले बाड़ें।
पहिले हमनी के ई देखब जा :

%%K t06-tit-1 sub
\VUpk{10.2pt}{40}{सुते के कमरा।}
\VUsaut{3.0mm}

%%K t06-c1-01-1 p
१. --- हमनी के गलत बेरा पहुँचल बानी जा, काहेकि कमरा के नौकरानी
\VUgras{कमरा} साफ करे में लागल बाड़ी।

%%K t06-c1-02-1 p
२. --- \VUgras{सिंगार-मेज} \textsuperscript{(1)} पर इहाँ-ओहिजा तरह-तरह
के चीज बा जवना से लोग हाथ-मुँह धोवेला, बाल झारेला — सार में सिंगार
करेला। ऊ ई बाड़ी सँ : \VUgras{चिलमची} \textsuperscript{(2)} आ
\VUgras{लोटा} \textsuperscript{(3)}, \VUgras{बाल के बुरुस}
\textsuperscript{(4)}, \VUgras{कंघी} \textsuperscript{(5)},
\VUgras{साबुन} \textsuperscript{(6)} आ \VUgras{स्पंज} \textsuperscript{(7)},
आ आखिर में दाँत के तरल दवाई खातिर \VUgras{सीसी} \textsuperscript{(8)}।

%%K t06-c1-03-1 p
३. --- भुइयाँ पर, सिंगार-मेज के आगे, ई रहल गंदा पानी जमा करे खातिर
\VUgras{बाल्टी} \textsuperscript{(9)}।

%%K t06-c1-04-1 p
४. --- \VUgras{अगला कमरा} \textsuperscript{(10)} में हमनी के कुछ
\VUgras{खूँटी} \textsuperscript{(13)} आ दू गो \VUgras{छाता}
\textsuperscript{(11)} देखत बानी जा, \VUgras{छाता-दान}
\textsuperscript{(12)} में।

%%K t06-c1-05-1 p
५. --- केवाड़ के दोसरा ओर \VUgras{सीसा वाला आलमारी}
\textsuperscript{(14)} बा जवना में \VUgras{कपड़ा-लत्ता}
\textsuperscript{(15)} रहेला।

%%K t06-c1-06-1 p
६. --- \VUgras{कमरा के नौकरानी} \textsuperscript{(6)} जब तक
\VUgras{पलंग} \textsuperscript{(17)} ठीक करत बाड़ी, तब तक ओकर अलग-अलग
हिस्सा देख लिहल जाव। \VUgras{पलंग के ढाँचा} \textsuperscript{(20)} में
नीक \VUgras{स्प्रिंग के गद्दा} \textsuperscript{(21)} बा, जवना पर
इउस्तीना अभिए \VUgras{ऊन के गद्दा} \textsuperscript{(22)} धरिहें। ओकरा
बाद ऊ कुरसी पर से दू गो \VUgras{चादर} \textsuperscript{(23)} आ
\VUgras{कंबल} \textsuperscript{(24)} लीहें। फेरु ऊ पलंग के सिरहाना पर
\VUgras{लमहर तकिया} \textsuperscript{(25)} आ \VUgras{तकिया}
\textsuperscript{(26)} धरिहें, जवन \VUgras{रजाई} \textsuperscript{(27)} के
साथे \VUgras{पलंग के दरी} \textsuperscript{(32)} पर बा। ऊ \VUgras{परदा}
\textsuperscript{(19)} आ \VUgras{पलंग के छत्तर} \textsuperscript{(18)} के
धूर झाड़े खातिर पंख के झाड़न चलावेली, आ लऽ, काम के एगो हिस्सा पूरा भइल।

%%K t06-c1-07-1 p
७. --- ओकरा बाद ओकरा \VUgras{पलंग के छोट मेज} \textsuperscript{(28)}
थोड़ा ठीक करे के परी, \VUgras{अलार्म घड़ी} \textsuperscript{(29)} में चाभी
भरे के परी, \VUgras{रात के दिया} \textsuperscript{(30)} तइयार करे के परी,
आ मोमबत्तीदान साफ करे खातिर \VUgras{मोमबत्ती} \textsuperscript{(31)}
निकाले के परी।

%%K t06-tit-2 sub
\VUsaut{5.0mm}
\VUpk{10.2pt}{40}{सिलाई के टोकरी आ अँगीठी के सामान।}
\VUsaut{3.0mm}

%%K t06-c1-08-1 p
८. --- \VUgras{सिलाई के टोकरी} \textsuperscript{(34)}, जे \VUgras{मोढ़ा}
\textsuperscript{(33)} लगे बा, ओकरो के ठीक करे के बहुते जरूरत बा। ओकरा
\VUgras{कढ़ाई} \textsuperscript{(35)} तह करे के परी, \VUgras{सिलाई के मेज}
\textsuperscript{(38)} में \VUgras{अँगुश्ताना}, \VUgras{धागा}
\textsuperscript{(36)}, \VUgras{सूई} \textsuperscript{(37)} आ
\VUgras{कैंची} \textsuperscript{(39)} बंद करे के परी, आ मेज के ढक्कन
\VUgras{चाभी} \textsuperscript{(40)} से बंद करे के परी।

%%K t06-c1-09-1 p
९. --- बीच के \VUgras{एक-गोड़ मेज} \textsuperscript{(41)} पर इउस्तीना
\VUgras{सुराही} \textsuperscript{(42)} के पानी बदल दीहें। ऊ
\VUgras{चीनीदान} \textsuperscript{(43)} में \VUgras{चीनी} धरिहें,
\VUgras{चीनी के चिमटी} \textsuperscript{(44)} साफ करिहें आ \VUgras{कपड़ा
के बुरुस} \textsuperscript{(46)} हटा दीहें।

%%K t06-c1-10-1 p
१०. --- कमरा के दाहिना ओर ओकरा से कम काम माँगी, काहेकि \VUgras{आराम-कुरसी}
\textsuperscript{(47)} आ ओकर \VUgras{कुसन} \textsuperscript{(48)} आपन ठीक
जगहा पर बा, आ काहेकि ठंढ नइखे, \VUgras{अँगीठी} \textsuperscript{(49)} के
सामान में से कुछ ना खिसकल बा। \VUgras{बेलचा} \textsuperscript{(50)},
\VUgras{चिमटा} \textsuperscript{(51)}, \VUgras{अँगीठी के टेक}
\textsuperscript{(55)}, \VUgras{राख-रोक} \textsuperscript{(56)} साफ बा ;
\VUgras{आग के परदा} \textsuperscript{(54)} अपना जगहा पर बा आ \VUgras{काठ
के संदूक} \textsuperscript{(52)} में \VUgras{लकड़ी के गुल्ला}
\textsuperscript{(53)} भरल बा।

%%K t06-noto-1 noto
\VUnotes{143.2mm}{%
(*) Babo = छोट बच्चा, अंगरेजी baby, फ्रेंच bébé, इतालवी bambino.%
}

%%K t06-noto-2 noto
\VUnotes{143.2mm}{%
(*) Lambrequino = फ्रेंच lambrequin, स्पेनी lambrequines, पुर्तगाली
lambrequins, जर्मन आ अंगरेजी में भी lambrequins, माने जर्मन, अंगरेजी,
फ्रेंच, पुर्तगाली, स्पेनी।%
}

%%K t06-c1-11-1 p
११. --- तबो ओकरा \VUgras{पेंडुलम घड़ी} \textsuperscript{(57)},
\VUgras{दीवट} \textsuperscript{(58)} आ \VUgras{जेब-दान}
\textsuperscript{(59)} के धूर झाड़े के परी, आ आखिर में \VUgras{सीसा}
\textsuperscript{(60)} पोंछे के परी।

%%K t06-c1-12-1 p
१२. --- रहल छोट बच्चा (बबुआ (*)) के \VUgras{पलना} \textsuperscript{(61)},
ऊ त पहिलही से तइयार बा।

%%K t06-tit-3 sub
\VUsaut{5.0mm}
\VUpk{10.2pt}{40}{बइठका।}
\VUsaut{3.0mm}

%%K t06-c2-01-1 p
१. --- अबहीं ई रहल बइठका। ई बहुते सुंदर \VUgras{कमरा} बा। जे अँगीठी दाहिना
कोना में बा ऊ नाजुक \VUgras{छोट मूरत} \textsuperscript{(1)} आ दू गो सुंदर
\VUgras{फूलदान} \textsuperscript{(3)} से सजल बा आ मखमली \VUgras{झालर}
\textsuperscript{(2)} (*) से ढँकल बा।

%%K t06-c2-02-1 p
२. --- चूल्हा के चिनगारी से कालीन ना जर जाव, एहीसे चूल्हा के आगे ताँबा
के \VUgras{चिनगारी-रोक} \textsuperscript{(4)} धइल गइल बा।

%%K t06-c2-03-1 p
३. --- ओकरा लगे एगो सुंदर \VUgras{लिखे के मेज} \textsuperscript{(5)} खुलल
बा। ओकरा पर \VUgras{घर के मालकिन} \textsuperscript{(23)} आपन चिट्ठी लिखत
बाड़ी।

%%K t06-c2-04-1 p
४. --- खिड़की पर \VUgras{खिड़की के परदा} \textsuperscript{(6)} लागल बा,
जवना के \VUgras{फीता} \textsuperscript{(8)} \VUgras{खूँटी}
\textsuperscript{(7)} में अँटकल बा, आ कपड़ा के बड़का परदा
\VUgras{ऊपर से लटकल} \textsuperscript{(9)} बा।

%%K t06-c2-05-1 p
५. --- खिड़की के ताख में \VUgras{गमला-खोल} \textsuperscript{(11)} में
हरियर \VUgras{पौधा} \textsuperscript{(10)} बा, एगो सुंदर खजूर जेकर पत्ता
लगभग \VUgras{मिट्टी तेल के लालटेन} \textsuperscript{(12)} तक फइलल बा।

%%K t06-c2-06-1 p
६. --- लालटेन के \VUgras{छत्ता} \textsuperscript{(13)} मारिया सजवले
बाड़ी, ऊ मनमोहक जवान लइकी \textsuperscript{(14)} जे पियानो
\textsuperscript{(15)} लगे बाड़ी। \VUgras{मोढ़ा} \textsuperscript{(16)} पर
सोझ बइठ के ऊ संगीत सीखत बाड़ी। ऊ बहुते होसियार आ सलीकेदार बाड़ी, जइसन
ओकर \VUgras{संगीत के आलमारी} \textsuperscript{(17)} से पता चलेला, जवन पूरा
सजल-सँवरल बा।

%%K t06-tit-4 sub
\VUsaut{5.0mm}
\VUpk{10.2pt}{40}{सैतान लइका।}
\VUsaut{3.0mm}

%%K t06-c2-07-1 p
७. --- ओकर भाई, पाउलुस, \VUgras{किताब} \textsuperscript{(19)} खोजत बा
\VUgras{किताब के आलमारी} \textsuperscript{(18)} में। ऊ \VUgras{किसोर}
\textsuperscript{(20)} बा, चंचल आ ना सहे लायक। जवन किताब बहुते ऊपर बा
ओकरा तक पहुँचे खातिर आलमारी में लटक के ऊ ओह \VUgras{मूरत}
\textsuperscript{(21)} के गिरावे के खतरा बनवत बा जवन ऊपर धइल बा।

%%K t06-c2-08-1 p
८. --- ऊ एह सैतानी खातिर एह बात के फायदा उठावत बा कि ओकर माई आपन
सहेली, एगो \VUgras{मेहरारू} \textsuperscript{(22)} जे कुछ दिन ओकरा घरे
बितावे आइल बाड़ी, से बतियावे में लागल बाड़ी। माई \VUgras{सोफा}
\textsuperscript{(24)} पर \VUgras{आरामकुरसी} \textsuperscript{(25)} लगे
बइठल बाड़ी, आ ओकर सहेली ओह \VUgras{कुरसी} \textsuperscript{(27)} पर बाड़ी
जेकर खाली \VUgras{पीठ के टेक} \textsuperscript{(26)} लउकत बा।

%%K t06-c2-09-1 p
९. --- भीत पर टँगल बड़का \VUgras{फ्रेम} \textsuperscript{(30)} में एगो
रिस्तेदारिन के \VUgras{तसवीर} \textsuperscript{(29)} बा। मेज पर धइल
\VUgras{एलबम} \textsuperscript{(32)} में बाकी घर के लोग के \VUgras{फोटो}
\textsuperscript{(33)} जुटल बा। लइकन अभिए ओकरा के उलटले बाड़ें, काहेकि
\VUgras{कुरसी} \textsuperscript{(31)} अबहीं ले मेज के चारो ओर जमा बा।

%%K t06-c2-10-1 p
१०. --- चीनी मिट्टी के \VUgras{चीनी फूलदान} \textsuperscript{(34)}, जे
\VUgras{कागज-छूरी} \textsuperscript{(35)} लगे बा, मारिया के ओकर परब पर
भेंट में मिलल रहे, आ जे \VUgras{आड़-परदा} \textsuperscript{(36)} कमरा के
कोना सजावत बा ऊ ओकर चाचा चीन से ले आइल रहलें।

%%K t06-c2-11-1 p
११. --- सुंदर \VUgras{चित्र} \textsuperscript{(37)} से सजल, जवन
\VUgras{भीत} \textsuperscript{(42)} पर टँगल बा, \VUgras{झाड़-फानूस} \textsuperscript{(38)}
से अँजोर, जेकर \VUgras{बिजली के बत्ती} \textsuperscript{(39)} साँझ के
खुसी से चमकेली, ई बइठका बहुते नीक लागत बा। काठ के फर्स पर बिछल मोलायम
\VUgras{कालीन} \textsuperscript{(40)}, आ ओतना आसान \VUgras{बिजली के घंटी}
\textsuperscript{(41)}, ओकर आराम पूरा कर देला।

%%K t06-tit-5 sub
\VUsaut{3.6mm}
\VUpk{10.2pt}{40}{खाए के कमरा।}
\VUsaut{3.0mm}

%%K t06-c3-01-1 p
१. --- खाना के घंट बाजल। मारिया के छोड़ के पूरा परिवार मेज के चारो ओर
जुटल बा। खाए के कमरा चीनी मिट्टी के बढ़िया \VUgras{अँगीठी}
\textsuperscript{(1)} से, जेकर \VUgras{नली} \textsuperscript{(2)} पातर बा,
नीक से गरम बा।

%%K t06-c3-02-1 p
२. --- एह कमरा में जादे सामान नइखे। भीत के किनारे, \VUgras{भीत-मेज}
\textsuperscript{(4)} के ऊपर, सजावटी \VUgras{छोट फुहारा}
\textsuperscript{(3)} टँगल बा। एकदम लगे \VUgras{परोसे के मेज}
\textsuperscript{(5)} बा जवना पर ठंढा खाना, मिठाई के थारी आ ऊ सब सामान
धइल जाला जवन तुरते ना चाहीं।

%%K t06-c3-03-1 p
३. --- एह तरह हमनी के ऊपर के तखता पर \VUgras{भुँजल मुरगी}
\textsuperscript{(7)}, \VUgras{मछरी} \textsuperscript{(8)} आ \VUgras{कटोरा}
\textsuperscript{(10)} देखत बानी जा जवना में \VUgras{फर}
\textsuperscript{(9)} बा। नीचे ई रहल \VUgras{पनीर} \textsuperscript{(11)}, \VUgras{डबलरोटी}
\textsuperscript{(12)}, आ \VUgras{तेल-सिरका के दान} \textsuperscript{(13)}
जवना में सलाद के छौंके के सब सामान बा : तेल, सिरका, \VUgras{नून}
\textsuperscript{(14)} आ \VUgras{गोलमिरचा} \textsuperscript{(15)}। आखिर
में, नीचला तखता पर \VUgras{केक} \textsuperscript{(17)} बा,
\VUgras{सरसों-दान} \textsuperscript{(16)} लगे।

%%K t06-c3-04-1 p
४. --- खाए के कमरा के घड़ी, जेकरा \VUgras{भीत के घड़ी}
\textsuperscript{(20)} कहल जाला, बहुते अलगे रूप के बा। ऊ काठ के फ्रेम में
जड़ल बा जवना में \VUgras{बैरोमीटर} \textsuperscript{(19)} आ
\VUgras{थर्मामीटर} \textsuperscript{(18)} भी बा।

%%K t06-tit-6 sub
\VUsaut{4.6mm}
\VUpk{10.2pt}{40}{खाना के परोसाई।}
\VUsaut{3.0mm}

%%K t06-c3-05-1 p
५. --- कमरा के सब भीत पर मोम लागल सागवान के \VUgras{काठ के पल्ला}
\textsuperscript{(21)} जड़ल बा। कपड़ा के \VUgras{केवाड़ के परदा}
\textsuperscript{(22)} ओह केवाड़ के ठीक से बंद करेला जवन बरामदा में
खुलेला। ओहिजा परिवार खाना के बाद कॉफी पिए जाई। \VUgras{प्याली}
\textsuperscript{(24)} आ \VUgras{तस्तरी} \textsuperscript{(25)} पहिलही से
\VUgras{थारी-मेज} \textsuperscript{(23)} पर तइयार बा।

%%K t06-c3-06-1 p
६. --- मेज के ऊपर छत से \VUgras{लटकन दिया} \textsuperscript{(26)} बान्हल
बा जेकर बहुते चमकत बत्ती साँझ के पूरा खाए के कमरा के अँजोर करेला।

%%K t06-c3-07-1 p
७. --- अबहीं \VUgras{खाए के मेज} \textsuperscript{(27)} के देख लिहल जाव। जब
परिवार आइल, खाना के सामान तइयार रहे। \VUgras{मेजपोस} \textsuperscript{(28)}
पर हर खाए वाला के जगहा पर \VUgras{थारी} \textsuperscript{(33)},
\VUgras{नैपकिन} \textsuperscript{(29)}, \VUgras{चम्मच} \textsuperscript{(34)},
\VUgras{काँटा} \textsuperscript{(35)}, \VUgras{छूरी} \textsuperscript{(36)}
आ \VUgras{गिलास} \textsuperscript{(32)} धइल रहे। सराब खातिर
\VUgras{बोतल} \textsuperscript{(30)} आ पानी खातिर \VUgras{सुराही}
\textsuperscript{(31)} भी रहे।

%%K t06-c3-08-1 p
८. --- \VUgras{नौकर} \textsuperscript{(39)} पहिले \VUgras{सूप के कटोरा}
\textsuperscript{(37)} ले अइलें, जवना में अबहीं ले थोड़ा सूप भाप छोड़त बा।
अबहीं ऊ पहिलका \VUgras{पकवान} \textsuperscript{(38)}, गोस के पकवान, परोसत
बाड़ें। ओकरा बाद ऊ ऊ सब देंहें जवन हमनी के गिना चुकल बानी जा, आ आखिर
में, मिठाई के बाद, ऊ एह बात के फायदा उठइहें कि ओकर मालिक लोग बरामदा में
चल गइल बा, आ मेज के सामान उठा के खाए के कमरा फेरु ठीक कर दीहें।

%%K t06-c3-08-2 p
\textit{नोट।} --- \textit{खाना-पीना से जुड़ल आम सबद इहाँ बहुते
संछेप में देहल गइल बा। ई सूची « होटल » (नं. १३) आ « बजार » (नं. १५)
वाला दू गो नकसा पर पूरा कइल जाई।}

%%K t06-tit-7 sub
\VUsaut{4.6mm}
\VUpk{10.2pt}{40}{रसोई।}
\VUsaut{3.0mm}

%%K t06-c4-01-1 p
१. --- रसोई में, जेकरा के हमनी के आधा खुलल केवाड़ से देखे लागत बानी जा,
एगो देखे लायक अस्त-बेस्त बा। \VUgras{चूल्हा} \textsuperscript{(1)} के
आगे, जवना में \VUgras{आग} \textsuperscript{(3)} चटकत बा, \VUgras{सीख
घुमावे वाला} \textsuperscript{(5)} लागल बा। सुंदर \VUgras{भुँजल गोस}
\textsuperscript{(7)} \VUgras{सीख में लागल} \textsuperscript{(6)} बा आ पकत
बा।

%%K t06-c4-02-1 p
२. --- \VUgras{धौंकनी} \textsuperscript{(8)} आ \VUgras{कोयला के बेलचा}
\textsuperscript{(9)}, जवना के अभिए इस्तेमाल भइल, \VUgras{लकड़ी}
\textsuperscript{(10)} नियर भुइयाँ पर पड़ल रह गइल।

%%K t06-c4-03-1 p
३. --- अँगीठी के ऊपर सब ठीक बा : \VUgras{लालटेन} \textsuperscript{(11)},
\VUgras{दियासलाई के डिब्बा} \textsuperscript{(13)}, जवना में
\VUgras{दियासलाई} \textsuperscript{(12)} बा, \VUgras{दीवट}
\textsuperscript{(14)} आ \VUgras{मोमबत्तीदान} \textsuperscript{(15)} आपन
\VUgras{मोमबत्ती} \textsuperscript{(16)} के साथे आपन जगहा पर बा। बाकिर
बड़का \VUgras{कोयला के बाल्टी} \textsuperscript{(18)}, जवना में
\VUgras{पथरकोयला} \textsuperscript{(17)} भरल बा आ जेकरा के
\VUgras{रसोइनिन} \textsuperscript{(23)} अभिए तहखाना से भर के लइले बाड़ी,
रसोई के बीचे रह गइल बा।

%%K t06-c4-04-1 p
४. --- सचहूँ, ओह बेचारी के जल्दी करे के बा ताकि खाना ठीक बेरा पर तइयार
हो जाव। अबहीं ऊ \VUgras{रसोई-चूल्हा} \textsuperscript{(19)} लगे बाड़ी आ
\VUgras{रसा} \textsuperscript{(24)} चलावत बाड़ी, जवन जादे जरे के ना
चाहीं।

%%K t06-c4-05-1 p
५. --- \VUgras{भट्ठी} \textsuperscript{(20)} में केक सिंकत बा, जवन ऊ
लइकन के जलपान खातिर बनवले बाड़ी। रसोई-चूल्हा के ऊपर \VUgras{करछुल}
\textsuperscript{(21)} आ \VUgras{झँझरी करछुल} \textsuperscript{(22)} टँगल
बा।

%%K t06-tit-8 sub
\VUsaut{4.0mm}
\VUpk{10.2pt}{40}{बरतन।}
\VUsaut{3.0mm}

%%K t06-c4-06-1 p
६. --- कमरा के छोर पर टँगल \VUgras{रसोई के बरतन} \textsuperscript{(25)}
बहुते साफ बा। ताँबा के सुंदर \VUgras{कड़ाही} \textsuperscript{(26)},
\VUgras{दूध के लोटा} \textsuperscript{(27)}, \VUgras{छननी}
\textsuperscript{(28)} आ \VUgras{कद्दूकस} \textsuperscript{(29)} जतन से
माँजल गइल बा।

%%K t06-c4-07-1 p
७. --- \VUgras{नून के डिब्बा} \textsuperscript{(30)} बरतन के आलमारी पर
\VUgras{चाय के केतली} \textsuperscript{(31)} आ \VUgras{तेल के बोतल}
\textsuperscript{(32)} लगे धइल बा। \VUgras{दराज} \textsuperscript{(33)}
में सायद खाए के सामान आ रसोई के छूरी रहेला।

%%K t06-c4-08-1 p
८. --- \VUgras{झाड़ू} \textsuperscript{(34)} आ \VUgras{पंख के झाड़न}
\textsuperscript{(35)} एगो कोना में बा। रसोइनिन अभिए \VUgras{काठ के
गुटका} \textsuperscript{(36)} इस्तेमाल कइले बाड़ी, ऊ ओकरा के खिड़की लगे
छोड़ देले बाड़ी।

%%K t06-c4-09-1 p
९. --- \VUgras{हाथ पोंछे के गमछा} \textsuperscript{(37)} भीत पर,
\VUgras{कीप} \textsuperscript{(38)} लगे टँगल बा। रसोई के ओह हिस्सा में
\VUgras{सिकाई के जाली} \textsuperscript{(39)}, \VUgras{काटे के छूरी}
\textsuperscript{(40)} आ \VUgras{काटे के पटरी} \textsuperscript{(41)} धइल
बा। ओकरा बाद, काटे के छूरी के ऊपर, एगो \VUgras{तखता} \textsuperscript{(42)}
पर \VUgras{बरतन} \textsuperscript{(43)}, \VUgras{देगची} \textsuperscript{(44)}
आपन \VUgras{ढक्कन} \textsuperscript{(45)} के साथे आ \VUgras{देग}
\textsuperscript{(46)} बा। \VUgras{तवा} \textsuperscript{(47)} लगहीं टँगल
बा।

%%K t06-c4-10-1 p
१०. --- एह कोना में खाली ताँबा के \VUgras{टोंटी} \textsuperscript{(48)}
चमकत बा। \VUgras{चौका} \textsuperscript{(49)}, जवना पर मोटका
\VUgras{मटका} \textsuperscript{(50)} धइल बा, आपन छोट आलमारी के साथे सायद
बहुते सुविधा वाला बा ; ओह आलमारी में \VUgras{सिरका के पीपा}
\textsuperscript{(51)} आपन \VUgras{टोंटी} \textsuperscript{(52)} के साथे,
\VUgras{कूड़ा के डिब्बा} \textsuperscript{(53)} आ \VUgras{मैल थारी}
\textsuperscript{(54)} रखल जाला।

%%K t06-tit-9 sub
\VUsaut{4.0mm}
\VUpk{10.2pt}{40}{रसोई में धइल दोसर चीज।}
\VUsaut{3.0mm}

%%K t06-c4-11-1 p
११. --- \VUgras{टब} \textsuperscript{(55)}, जवना में \VUgras{तरकारी}
\textsuperscript{(58)} धोवल जाला, आ \VUgras{लोहा के कड़ाह}
\textsuperscript{(56)}, जवन \VUgras{फर्स के टाइल} \textsuperscript{(57)} पर
धइल बा, सायद ओहिजे ओह आलमारी में रहेला।

%%K t06-c4-12-1 p
१२. --- रसोइनिन के चूल्हा के कमी त पक्का नइखे, काहेकि ई रहल मेज के
कोना पर \VUgras{गैस के चूल्हा} \textsuperscript{(59)} जवना पर छोट कड़ाही
में पानी गरम होत बा। जवन \VUgras{भाप} \textsuperscript{(60)} ओहसे निकलत
बा ऊ बतावत बा कि सायद ऊ उबलत बा। सायद ई पानी कॉफी खातिर काम आई, काहेकि
ई रहल मेज के दोसरा छोर पर \VUgras{कॉफी के केतली} \textsuperscript{(64)}
आ \VUgras{कॉफी के चक्की} \textsuperscript{(63)}।

%%K t06-c4-13-1 p
१३. --- चूल्हा \VUgras{गैस के मुँह} \textsuperscript{(62)} से लमहर
\VUgras{रबर के नली} \textsuperscript{(61)} के जरिए जुड़ल बा।

%%K t06-c4-14-1 p
१४. --- \VUgras{दिहाड़ी के नौकरानी} \textsuperscript{(66)}, जे रसोइनिन के
मदद करे आवेली, खाना खातिर तरकारी तइयार करत बाड़ी। जब ऊ छिला आ धोवा
जाई, त ऊ ओकरा के \VUgras{परात} \textsuperscript{(65)} में धर दीहें आ
पकावे के बेरा के अगोरिहें।

%%K t06-tit-10 sub
\VUsaut{4.0mm}
{\centering\fontsize{10.2pt}{10.2pt}\selectfont\textls[40]{\scshape नहाए के कमरा।} \textit{(गुसलखाना।)}\par}
\VUsaut{3.0mm}

%%K t06-c5-01-1 p
१. --- घर के खास कमरा जाने खातिर हमनी के खाली एगो ढिठाई करे के बाकी
बा : नहाए के कमरा के बनावट देखल। एहमें जादे बेरा ना लागी, काहेकि ऊ
बड़ नइखे, आ हमनी के जल्दिए जान लेब जा कि \VUgras{नहाए के टब}
\textsuperscript{(1)} में नीक \VUgras{पानी गरम करे के भट्ठी}
\textsuperscript{(2)} बा, कि \VUgras{साबुनदानी} \textsuperscript{(3)}
नहाए वाला के हाथ तक पहुँचेला, आ कि \VUgras{गमछा-टाँग} \textsuperscript{(5)}
से \VUgras{गमछा} \textsuperscript{(4)} पक्का आसानी से सूख जाई।

%%K t06-c5-02-1 p
२. --- एह कमरा के भीत \VUgras{चीनी मिट्टी के टाइल} \textsuperscript{(6)}
से ढँकल बा, जवना के चलते \VUgras{फुहारा} \textsuperscript{(7)} लगावल संभव
भइल बिना एह डर के कि ओकरा इस्तेमाल में जवन पानी छिटकेला ऊ भीत के बिगाड़
दी। जस्ता के \VUgras{छोट परात} \textsuperscript{(8)}, जवन गोड़ धोवे के
काम आवेला, एह सामान के पूरा कर देला।
